% ============================================================
% CHAPTER 1 — Introduction (REVISED v2)
% ============================================================

\chapter{Introduction}
\label{chap:introduction}

\ac{BPM} systems are increasingly required to operate in dynamic and
flexible environments where processes cannot always be represented as
rigid execution flows \cite{Reichert2012EnablingTechnologies}.
Declarative process modelling approaches address this challenge by
specifying the constraints that process executions must satisfy, instead
of prescribing a single fixed sequence of activities \cite{PesicDECLARE:Processes}.
Among these approaches, \ac{DCR} graphs have gained attention due to
their formal semantics and their ability to represent flexible yet
compliant business processes \cite{Hildebrandt2011DeclarativeGraphs}.

While \ac{DCR} graphs guarantee that process executions satisfy the
specified constraints, they provide limited support for identifying
which compliant executions are operationally preferable, a problem
characterised in the literature as the exploration of a flat conformance
space \cite{AndreaBurattin2021Exploring_Conformance_Space}.

In practice, a process may contain many valid traces with very different
characteristics in terms of execution cost, duration, or resource
consumption \cite{Diaz2024Pareto-OptimalModels}.
Optimising process execution within this large conformance space
therefore becomes an important challenge.

\ac{RL} offers a promising approach for process optimisation because
agents can learn execution strategies directly through interaction with
an environment \cite{Sutton2018ReinforcementIntroduction}.
However, standard \ac{RL} approaches do not inherently guarantee process
compliance. An agent may learn behaviours that maximise reward while
violating critical business constraints
\cite{Garcia2015ALearning}, which makes purely neural
optimisation problematic for compliance-sensitive domains.

This thesis proposes a compliance-aware multi-objective optimisation
framework that combines \ac{RL} with \ac{DCR} graphs.
The approach models \ac{DCR} execution as a sequential decision-making
problem where the agent learns to optimise process execution while
respecting the semantics of the underlying declarative model.
Compliance is enforced through a shielding mechanism
\cite{Alshiekh2017SafeShielding} that prevents invalid actions during both training
and execution, while a multi-objective reward function guides the agent
towards Pareto-optimal traces according to business objectives such as
cost and duration.
The objectives and research questions guiding this investigation are
stated in Section~\ref{sec:research_questions}, and the specific
contributions of this work are presented in
Section~\ref{sec:contributions}.

% ----------------------------------------------------------------
\section{Problem Overview}
\label{sec:problem_overview}

The flexibility of declarative process models comes at the cost of a
large and weakly differentiated conformance space, in which existing
\ac{DCR}-based systems can verify that traces are valid but cannot rank
them by operational quality.
Current optimisation approaches present complementary limitations.
Exact symbolic techniques guarantee correctness and optimality, but
their underlying decision problems scale poorly as process complexity increases \cite{Diaz2024Pareto-OptimalModels, Felli2023Data-awareSMT}.

Neural approaches such as \ac{RL} are more scalable and adaptive, but
they may generate traces that violate process constraints
\cite{Garcia2015ALearning}.

Consequently, there is currently a gap between:
\begin{itemize}
    \item optimisation approaches that guarantee compliance but struggle
    with scalability,
    \item and scalable learning-based approaches that lack formal
    compliance guarantees.
\end{itemize}

This gap, substantiated through a systematic comparison of related work
in Section~\ref{sec:research_gap},
motivates the central proposal of this thesis, an architecture in which
a learning agent explores the conformance space freely while a symbolic
shield derived from the \ac{DCR} semantics makes constraint violations
structurally impossible.

% ----------------------------------------------------------------
\section{Objectives and Research Questions}
\label{sec:research_questions}

The main objective of this thesis is to develop and evaluate a
compliance-aware \ac{RL} framework for multi-objective optimisation of
declarative business processes.
This involves modelling \ac{DCR} execution as a sequential
decision-making problem, designing a reward structure that captures
multiple business objectives simultaneously, enforcing compliance
through symbolic shielding, and evaluating both the external performance
of the resulting framework against exact and heuristic baselines and the
internal contribution of each of its components.

\paragraph{Main Research Question (MRQ)}
How can \ac{RL} optimise declarative business-process execution while
guaranteeing compliance with \ac{DCR} constraints?

\paragraph{Sub-Research Questions (SRQs)}
\begin{enumerate}
    \item How can \ac{DCR} graph semantics be represented within an
    \ac{RL} environment?

    \item To what extent can \ac{PPO}-based agents discover
    Pareto-optimal traces in declarative process models?

    \item How does reward design influence convergence and compliant
    agent behaviour?

    \item Can the framework learn resource-allocation strategies through
    role-conditioned action spaces?
\end{enumerate}

These research questions are formulated at the level of knowledge
claims. Chapter~\ref{chap:experimental_setup} operationalises them as
five concrete, measurable evaluation questions (EQ1 through EQ5) and
makes the mapping between research questions, evaluation questions,
contributions, and evidence chapters explicit.

% ----------------------------------------------------------------
\section{Scope and Limitations}
\label{sec:scope}

This thesis focuses on compliance-aware, multi-objective optimisation of
declarative business processes represented as \ac{DCR} graphs, covering
the optimisation of cost and duration objectives and resource-aware
execution through role-conditioned actions.

To maintain a feasible scope, several aspects are excluded or
simplified:
\begin{itemize}
    \item The framework focuses exclusively on \ac{DCR} graphs and does
    not address other declarative languages in detail,
    \item The \ac{RL} implementation is limited to \ac{PPO}-based agents,
    \item The Pareto front is approximated through sampling rather than
    formally guaranteed,
    \item The empirical validation is limited to benchmark processes,
    two used for development (Expense Report and Loan Application) and a
    suite of independently mined \ac{DCR} graphs used for external
    validation (Chapter~\ref{chap:external_validation}), rather than a
    deployment in a live production environment,
    \item Large-scale industrial deployment and distributed execution are
    outside the scope,
    \item Partial observability, online learning, and
    continuous-time process execution are not considered.
\end{itemize}

% ----------------------------------------------------------------
\section{Contributions}
\label{sec:contributions}

This thesis makes three primary contributions to the field of
compliance-aware business process optimisation.

\textbf{C1.\label{cont:c1} A DCR-shielded reinforcement learning environment.}
A reinforcement learning environment is introduced that tightly couples
\ac{DCR} graph semantics with a shielding mechanism intrinsic to the
\ac{DCR} execution engine.
In contrast to approaches that discourage non-compliant behaviour
through penalty terms or Lagrangian relaxation
\cite{AltmanCONSTRAINEDPROCESSES, Achiam2017ConstrainedOptimization}, where
constraint violations still occur during training, the agent here
samples over the full action space and any non-enabled action is
intercepted by the engine before it can take effect, leaving the
process state entirely unchanged.
Compliance is therefore a hard structural property of every trace
generated throughout training, rather than a statistical tendency that
emerges only after convergence.
To the best of our knowledge, this constitutes the first application of
shielded reinforcement learning \cite{Alshiekh2017SafeShielding} to \ac{DCR} graph
semantics.

\textbf{C2.\label{cont:c2} A multi-objective scalarised reward framework for Pareto-optimal trace generation.}
A normalised reward function is proposed that decomposes into
structural, cost, and duration components, enabling independent \ac{PPO}
agents trained under different scalarisation weights $(\alpha, \beta)$
\cite{DasDr.Das-NBIPaper-1998}
to approximate the Pareto front of valid process executions.
On the Expense Report benchmark, the framework recovers the exact Pareto
front identified by a CSP+COP exact solver \cite{Diaz2024Pareto-OptimalModels}.
On the Loan Application benchmark, it produces a well-populated front of
26 non-dominated points while maintaining full compliance throughout
training.

\textbf{C3.\label{cont:c3} A role-conditioned extension of \ac{EDCR} for
resource-aware optimisation.}
The base framework is extended to $(\text{event}, \text{role})$ action
pairs, so the agent optimises cost and duration at the level of
individual activities by choosing which resource role executes each
event. Per-role cost and duration multipliers, structured after the
resource profiles of Kirchdorfer et
al.~\cite{KirchdorferLukasFromPolicies}, are encoded directly in the
\ac{DCR} XML specification, preserving compatibility with existing
\ac{EDCR} graphs. Under unlimited role availability this choice is
separable across events; to turn it into a genuine allocation decision,
the framework is further extended with a per-case \emph{Expert capacity
budget}, enforced by the same shield, which couples the per-event
choices into an opportunity-cost allocation problem. Experiments on the Loan Application process show that, under unlimited role availability, the policy recovers the closed-form optimal role assignment, a separable pattern that distinguishes objective-sensitive
from role-indifferent activities and is consistent with the cost--duration trade-off structure independently identified by evolutionary and simulation-based
approaches~\cite{KirchdorferLukasFromPolicies,
ArtemisDoumeniParetoAgentSimulator}. Under the Expert capacity budget, where the per-event choices are no longer separable, the agent learns a genuine allocation policy that outperforms standard allocation
heuristics under scarcity.

Taken together, C1--C3 constitute the first unified framework combining
declarative \ac{DCR} semantics, compliance-by-construction shielding,
multi-objective optimisation, and resource-aware decision-making within
a single reinforcement learning pipeline.
While prior work addresses subsets of these properties in isolation, no
existing approach provides formal compliance guarantees, Pareto-optimal
trace generation, and event-level role-conditioned optimisation simultaneously
over a declarative process model.

% ----------------------------------------------------------------
\section{Structure of the Thesis}
\label{sec:structure_thesis}

The remainder of this thesis is organised as follows.

Chapter~\ref{chap:background} introduces the theoretical foundations
underlying the work, including declarative process modelling, \ac{DCR}
graphs, \ac{RL}, and multi-objective optimisation.

Chapter~\ref{chap:literarture} reviews related work on process
optimisation, compliance-aware \ac{RL}, declarative process execution,
and resource-aware optimisation, and derives the research gap addressed
by this thesis.

Chapter~\ref{chap:method} presents the proposed system architecture and
methodology, including the \ac{DCR} \ac{RL} environment, reward
formulation, shielding mechanism, \ac{PPO} training setup, and
role-conditioned extension.

Chapter~\ref{chap:experimental_setup} describes the experimental setup,
stating the evaluation questions, benchmarks, training protocol,
baselines, evaluation metrics, and reproducibility information.

Chapter~\ref{chap:results} presents the experimental results on both
benchmarks: the validation of Pareto-optimal trace discovery, the effect
of reward weighting, and the role-conditioned optimisation experiments.

Chapter~\ref{chap:internal_validation} provides an internal validation
of the methodology, isolating the individual contribution of the shield
and the learning agent through an ablation, and assessing the
generalisation of the learned policies across seeds and a temporal
holdout.

Chapter~\ref{chap:external_validation} tests external validity,
re-training the framework from scratch on a suite of independently mined
\ac{DCR} graphs never used during its development, and reports the
resulting compliance replication and scalability frontier.

Chapter~\ref{chap:discussion} discusses the implications, limitations,
and scalability considerations of the proposed framework, positions it
against prior approaches, and revisits the research questions.

Chapter~\ref{chap:conclusion} summarises the contributions, answers the
main research question, and outlines directions for future work.
